%%%%%%%%%%%%%%%%%%%%%%%%%%%%%%%%%%%%%%%%%
% Short Sectioned Assignment LaTeX Template Version 1.0 (5/5/12)
% This template has been downloaded from: http://www.LaTeXTemplates.com
% Original author:  Frits Wenneker (http://www.howtotex.com)
% License: CC BY-NC-SA 3.0 (http://creativecommons.org/licenses/by-nc-sa/3.0/)
%%%%%%%%%%%%%%%%%%%%%%%%%%%%%%%%%%%%%%%%%

% \documentclass[paper=a4, fontsize=11pt]{scrartcl} % A4 paper and 11pt font size
\documentclass[11pt, a4paper]{book}
\usepackage[T1]{fontenc} % Use 8-bit encoding that has 256 glyphs
\usepackage[utf8]{inputenc}
\usepackage{fourier} % Use the Adobe Utopia font for the document - comment this line to return to the LaTeX default
\usepackage{listings} % para insertar código con formato similar al editor
\usepackage[spanish, es-tabla]{babel} % Selecciona el español para palabras introducidas automáticamente, p.ej. "septiembre" en la fecha y especifica que se use la palabra Tabla en vez de Cuadro
\usepackage{url} % ,href} %para incluir URLs e hipervínculos dentro del texto (aunque hay que instalar href)
\usepackage{graphics,graphicx, float} %para incluir imágenes y colocarlas
\usepackage[gen]{eurosym} %para incluir el símbolo del euro
\usepackage{cite} %para incluir citas del archivo <nombre>.bib
\usepackage{enumerate}
\usepackage{hyperref}
\usepackage{graphicx}
\usepackage{tabularx}
\usepackage{booktabs}

%
% Para insertar código
%

\usepackage{listings}
\usepackage{color}
\definecolor{lightgray}{rgb}{.9,.9,.9}
\definecolor{darkgray}{rgb}{.4,.4,.4}
\definecolor{purple}{rgb}{0.65, 0.12, 0.82}

\lstdefinelanguage{JavaScript}{
  keywords={typeof, new, true, false, catch, function, return, null, catch, switch, var, if, in, while, do, else, case, break},
  keywordstyle=\color{blue}\bfseries,
  ndkeywords={class, export, boolean, throw, implements, import, this},
  ndkeywordstyle=\color{darkgray}\bfseries,
  identifierstyle=\color{black},
  sensitive=false,
  comment=[l]{//},
  morecomment=[s]{/*}{*/},
  commentstyle=\color{purple}\ttfamily,
  stringstyle=\color{red}\ttfamily,
  morestring=[b]',
  morestring=[b]"
}

\lstset{
   language=JavaScript,
   backgroundcolor=\color{lightgray},
   extendedchars=true,
   basicstyle=\footnotesize\ttfamily,
   showstringspaces=false,
   showspaces=false,
   numbers=left,
   numberstyle=\footnotesize,
   numbersep=9pt,
   tabsize=2,
   breaklines=true,
   showtabs=false,
   captionpos=b
}

%%%%%%%%%%%%%%%%%%%%%%%%%%%%%%%%%%%%

%problema fuentes tamaños
\usepackage{lmodern}

\usepackage[table,xcdraw]{xcolor}
\hypersetup{
	colorlinks=true,	% false: boxed links; true: colored links
	linkcolor=black,	% color of internal links
	urlcolor=cyan		% color of external links
}
\renewcommand{\familydefault}{\sfdefault}
\usepackage{fancyhdr} % Custom headers and footers
\pagestyle{fancyplain} % Makes all pages in the document conform to the custom headers and footers
\fancyhead[L]{} % Empty left header
\fancyhead[C]{} % Empty center header
\fancyhead[R]{Víctor González Argudo} % My name
\fancyfoot[L]{} % Empty left footer
\fancyfoot[C]{} % Empty center footer
\fancyfoot[R]{\thepage} % Page numbering for right footer
%\renewcommand{\headrulewidth}{0pt} % Remove header underlines
\renewcommand{\footrulewidth}{0pt} % Remove footer underlines
\setlength{\headheight}{13.6pt} % Customize the height of the header

\usepackage{titlesec, blindtext, color}
\definecolor{gray75}{gray}{0.75}
\newcommand{\hsp}{\hspace{20pt}}
\titleformat{\chapter}[hang]{\Huge\bfseries}{\thechapter\hsp\textcolor{gray75}{|}\hsp}{0pt}{\Huge\bfseries}
\setcounter{secnumdepth}{4}
\usepackage[Lenny]{fncychap}


\begin{document}

	% Plantilla portada UGR
	\begin{titlepage}
\newlength{\centeroffset}
\setlength{\centeroffset}{-0.5\oddsidemargin}
\addtolength{\centeroffset}{0.5\evensidemargin}
\thispagestyle{empty}

\noindent\hspace*{\centeroffset}\begin{minipage}{\textwidth}

\centering
\includegraphics[width=0.9\textwidth]{logos/logo_ugr.jpg}\\[1.4cm]

\textsc{ \Large TRABAJO FIN DE GRADO\\[0.2cm]}
\textsc{ GRADO EN INGENIERIA INFORMATICA}\\[1cm]

{\Huge\bfseries MemeHub \\}
\noindent\rule[-1ex]{\textwidth}{3pt}\\[3.5ex]
{\large\bfseries Edición Online de Imágenes }
\end{minipage}

\vspace{2.5cm}
\noindent\hspace*{\centeroffset}
\begin{minipage}{\textwidth}
\centering

\textbf{Autor}\\ {Víctor González Argudo}\\[2.5ex]
\textbf{Director}\\ {Juan Julián Merelo Guervós}\\[2cm]
\includegraphics[width=0.3\textwidth]{logos/etsiit_logo.png}\\[0.1cm]
\textsc{Escuela Técnica Superior de Ingenierías Informática y de Telecomunicación}\\
\textsc{---}\\
Granada, Junio de 2021
\end{minipage}
\end{titlepage}


	% Plantilla prefacio UGR
	\thispagestyle{empty}

\begin{center}
{\large\bfseries MemeHub \\ Edición Online de Imágenes }\\
\end{center}
\begin{center}
Víctor González Argudo\\
\end{center}

%\vspace{0.7cm}

\vspace{0.5cm}
\noindent{\textbf{Palabras clave}: \textit{Javascript, React.js, JSX, Editor de Imágenes, Software Libre}}
\vspace{0.7cm}

\noindent{\textbf{Resumen}\\}
	

\cleardoublepage

\begin{center}
	{\large\bfseries MemeHub \\ Online Image Editing}\\
\end{center}
\begin{center}
	Víctor González Argudo\\
\end{center}
\vspace{0.5cm}
\noindent{\textbf{Keywords}: \textit{Javascript, React.js, JSX, Image Editor, Open Source}}
\vspace{0.7cm}

\noindent{\textbf{Abstract}\\}


\cleardoublepage

\thispagestyle{empty}

\noindent\rule[-1ex]{\textwidth}{2pt}\\[4.5ex]

D. \textbf{Tutora/e(s)}, Profesor(a) del ...

\vspace{0.5cm}

\textbf{Informo:}

\vspace{0.5cm}

Que el presente trabajo, titulado \textit{\textbf{Chief}},
ha sido realizado bajo mi supervisión por \textbf{Estudiante}, y autorizo la defensa de dicho trabajo ante el tribunal
que corresponda.

\vspace{0.5cm}

Y para que conste, expiden y firman el presente informe en Granada a Junio de 2018.

\vspace{1cm}

\textbf{El/la director(a)/es: }

\vspace{5cm}

\noindent \textbf{(nombre completo tutor/a/es)}

\chapter*{Agradecimientos}






	% Índice de contenidos
	\newpage
	\tableofcontents

	% Índice de imágenes y tablas
	\newpage
	\listoffigures

	% Índice de código
	\newpage
	\renewcommand\lstlistingname{Código}
	\renewcommand\lstlistlistingname{Índice de código}
	\lstlistoflistings

	% Si hay suficientes se incluirá dicho índice
	\listoftables 
	\newpage

	% Introducción 
	\chapter{Introducción}

Este proyecto es software libre, y está liberado con la licencia \cite{gplv3}.

%//los memes blablabla y editores en la actualidad blablabla

%Por qué hay necesidad de mi apliación 

	% Descripción del problema y hasta donde se llega
	\chapter{Descripción del problema}

Este proyecto

//accesibiliddad blabla bla



	% Estado del arte
	% 	1. Crítica al estado del arte
	% 	2. Propuesta
	\chapter{Estado del arte}

La motivación de este proyecto surge precisamente por el estado actual de los editores de imágenes. 
La mayoría de usuarios, edita el contenido que sube online desde una aplicación local instalada en su ordenador (Microsoft Paint, Photoshop), o en el caso de los que emplean aplicaciones en la nube, descargan sus imágenes para luego resubirlas a sus redes.
Concretamente los memes suelen ser ediciones especialmente simples, donde simplemente se añade un texto, se superponen varias imágenes o se dibuja algo encima de una imagen base.

Este proyecto tiene como objetivo tratar este problema y evitar en la medida de lo posible el uso del almacenamiento local de imágenes, tanto para crear, como para subir el contenido ya editado. 


	
	% Análisis del problema
	% 1. Análisis de requisitos
	% 2. Análisis de las soluciones
	% 3. Solución propuesta
	% 4. Análisis de seguridad
	\chapter{Análisis del problema}
 
%//Estudio de viabilidad (tiene sentido)
%//Perfiles implicados, trabajadores, usuarios, 
%//que va a hacer el cliente (tendra exito si tiene una experiencia fluida)

%//Casos de uso
%//Req fun
%//Req no fun



	\chapter{Planificación}

\section{Metodología utilizada}
    Para la realización de este proyecto se ha seguido una metodología ágil o 'Agile', esta metodología
    es especialmente buena para los proyetos que necesitan de mucha flexibilidad, por ende, por 
    la naturaleza de algo como una aplicación de edición de imagen, que es algo que invita fácilmente a añadir nuevas funcionalidades
    y a cambiar o adaptar antiguas al nuevo contenido que se añada, ha sido escogida para el desarrollo de este trabajo.

    Para la metodología ágil escogida 
\section{Temporización}
    
\section{Seguimiento del desarrollo}
    %kanban

	% Desarrollo bajo sprints: 
	% 	1. Permitir registros y login de usuarios
	% 	2. Desarrollo del sistema de incidencias
	% 	3. Desarrollo del sistema de denuncias administrativas y accidentes
	% 	4. Desarrollo del sistema de croquis
	%   5. Instalación de la aplicación de manera automática
	\chapter{Implementación}

La implementación del software se ha dividido en hitos. Estos, han sido definidos en Github
y cada uno de ellos contiene un grupo de \textit{issues} que se corresponden con las distintas
mejoras que se han ido incorporando al software a lo largo de su desarrollo.\\

Esquema para luego redactar bien:

- Empezamos aprendiendo React

- Hice una demo básica

- Busqué librerias y cosas mientras seguí aprendiendo React

- Despues de unos cuantos proyectos de prueba y elegidas las librerias me lancé con el editor

- A partir de la demo rehice todo para adaptarlo a la libería 

- Lo primero fue hacer que se pudiera añadir una foto base (DragAndDrop y Konva canvas)

- Explicar todo tema lienzo, cálculos etc..

- Se añadió poder quitar la foto y volver a añadir otra

- Comenzamos el poder añadir texto

- todo tema texto, textarea, propiedades, drag, transformador, lo de adaptar el texto al tamaño de la caja bla bla 

- Comienzo de la toolbar para poder poner el texto bonito

- Tema toolbar, material-ui, como mandamos la info desde la toolbar al texto y lo actualizamos bla bla

- Hablar de ajustes varios, tema eventos etc.. clicks handlers bla bla

- Añadido el ctrl + v de imagen principal (esto quizá explicarlo después para juntar todo el tema texto)

- Handleo de varios textos diferentes

- Cuando ya se cambiaban los textos desde la toolbar ahora como se cambia la toolbar para que coincida con las opciones del texto

- Hablar un poco de tema componentes funcionales vs clases (todos los cambios que se hicieron)

- PRIMERA TANDA DE TESTS (preguntar a JJ si explicar esto todo luego o aquí)

- Además de texto ahora a añadir mas cosas, Imágenes y sus transformaciones

- Hablar del tema lineas, como se movian las cosas al pintar encima, la lógica de las herramientas
  como se ha solucionado, porque se ha hecho así y el React.clone()

- [pensar lo que falta por aqui]

- tema eventos, copiar pegar, guardar etc

	% Presupuesto

	% Conclusiones
	\chapter{Conclusiones y trabajos futuros}

Si estoy satisfecho

Has conseguido cumplir los objetivos

se ha seguido la planificacion



	% Trabajos futuros


	
	\newpage
	\bibliography{bibliografia}
	\bibliographystyle{plain}
	
\end{document}

